\documentclass[11pt]{article}
\usepackage[a4paper,margin=25mm]{geometry}
\usepackage{amsmath,amssymb,amsthm,booktabs,tikz}
\usepackage{xr}
\externaldocument[main-]{five-colour-cubes}
\usepackage[colorlinks=true,urlcolor=blue,citecolor=blue]{hyperref}
\hypersetup{pdftitle={Sharp thresholds for five-colour cube assemblies},pdfauthor={Jason Davies}}
\newtheorem{theorem}{Theorem}
\newtheorem{lemma}[theorem]{Lemma}
\newtheorem{proposition}[theorem]{Proposition}
\newtheorem{corollary}[theorem]{Corollary}
\theoremstyle{remark}
\newtheorem*{remark}{Remark}
\newcommand{\col}{\{0,1,2,3,4\}}
\newcommand{\mass}[2]{#1(#2)}
\title{Certificate proof and verification supplement}
\author{Jason Davies}
\date{November 4, 2026}
\begin{document}
\maketitle
This supplement preserves the computer-assisted corner proof used by the
Lean development. It is independent of the elementary corner argument in
the revised article and is not a formalisation of that argument.

\section{Reduction to finite problems}\label{sec:cases}
Eight copies of any variety solve that variety, so a solution-free collection
has every multiplicity at most seven. Suppose such a collection has size 15.

\begin{lemma}[Anchored Hall cover]\label{lem:anchor}
Suppose a collection contains $m$ copies of a variety $A$ and does not solve
$A$. There is a set $K$ of its target positions such that
\begin{equation}\label{eq:anchor}
 |K|\le 7-m,\qquad
 \#\{\text{other copies fitting some position outside }K\}
       \le 7-m-|K|.
\end{equation}
Conversely, such a set certifies failure of target $A$.
\end{lemma}
\begin{proof}
Apply~\eqref{main-eq:hall} to $A$ and put $K=D^c$. Every copy of $A$ is a
neighbour of the nonempty set $D$. Subtract these $m$ copies from the bound
$|D|-1=7-|K|$. The converse is the same Hall inequality read backwards.
\end{proof}

First suppose the maximum multiplicity $m$ is at least two. Choose an
adjacent-repeat anchor attaining $m$ if possible; otherwise choose an
opposite-repeat anchor attaining $m$. In the latter case all adjacent-repeat
multiplicities are at most $m-1$. Relabel the colours to send the anchor to
a fixed representative of its orbit.

For each cover $K$ in Lemma~\ref{lem:anchor}, put
\[
 S(K)=N_A(K^c)\setminus\{A\},\qquad b(K)=7-m-|K|.
\]
Write $S=S(K)$ and $b=b(K)$. The case has
$x_A=m$, $x(S)\le b$, and the multiplicity bounds just described.
Additional upper bounds use two-variety solutions. If $j$ copies of $A$
and $8-j$ copies of $V$ build a target, with $j\le m$, then
$x_V\le 7-j$ in a solution-free collection. These tests are finite:
the two varieties must cover all eight target positions, $A$ must fit at
least $j$, and $V$ at least $8-j$. They are also sufficient: at most $j$
positions are unavailable to $V$,
and all of them are available to $A$. Assign these to $A$, extend the
assignment to $j$ positions that $A$ fits, and assign the remaining
$8-j$ positions to $V$.

Colour permutations and simultaneous reflection of the entire collection
preserve solvability. First identify equal pairs $(S(K),b(K))$, even when
they arise from inequivalent position subsets $K$; then retain one
representative per orbit of these pairs under the anchor stabiliser.
The counts below refer to these constraint pairs. Transporting a cover
also transports its constrained varieties, budget and multiplicities, so
no counterexample is lost. Global reflection transports the whole
collection; individual cubes still move only by proper rotations.
The repeated-variety cases give 244 constraint orbits; the breakdown is in
Appendix~\ref{app:verification}.

For distinct collections a different choice of anchor reduces the count.
\begin{lemma}[Shared two-colour corner]\label{lem:shared}
Among eleven adjacent-repeat cubes, two share a two-colour corner type.
If a distinct collection contains such a pair and one member is chosen as
anchor, the set consisting of its two two-colour positions is not
Hall-deficient.
\end{lemma}
\begin{proof}
Each adjacent-repeat cube with repeated colour $r$ has exactly two such
types, $rrx$ and $rry$, with $x\ne y$. There are only $5\cdot4=20$ types $rrx$. Eleven cubes contribute
22 incidences, so two cubes share a type. For the two characteristic positions
of the chosen anchor, the anchor itself and the other cube are two distinct
neighbours. Their neighbourhood therefore has size at least two.
\end{proof}
If an opposite-repeat variety is present, choose it as anchor; this gives
31 constraint orbits with every multiplicity at most one. Otherwise use
Lemma~\ref{lem:shared} to choose an adjacent-repeat anchor. Omit the cover
complementary to its two characteristic positions and set all opposite-repeat
capacities to zero. Its stabiliser preserves these positions, so the omission
is compatible with the orbit reduction. It removes one of the 83
adjacent-repeat constraint orbits, leaving 82 further cases.

\begin{proposition}\label{prop:cover}
Every solution-free fifteen-cube collection is carried by colour relabelling
and, if needed, global reflection into one of the 357 exported root cases.
\end{proposition}
\begin{proof}
The preceding arguments cover repeated and distinct collections. The finite
catalogue, two-variety bounds and coverage of all anchor-cover orbits are
checked from physical rotations by the accompanying audit. It regenerates
rotations from quarter turns, separately from the signed-axis permutations
used by the case generator, and checks every exported bound and every cover
orbit. The resulting counts are
\[
  244+31+82=357.
\]
The cases need not be disjoint; completeness of the cover is what is required.
\end{proof}

Each root specifies an anchor $A$ of multiplicity $m$, coordinate bounds
$\ell,u$, a set $S$ of non-anchor varieties and a budget $b$.
The bounds fix $\ell_A=u_A=m$, with $\ell_i=0$ for $i\ne A$.
A vector $x\in\mathbb Z_{\ge0}^{75}$ is \emph{admissible} for this root if
\begin{equation}\label{eq:admissible}
\begin{gathered}
 \ell\le x\le u,\qquad \sum_{i=1}^{75}x_i=15,\qquad x(S)\le b,\\
 \text{for every target }T\text{, there is }\varnothing\ne D\subseteq\{1,\ldots,8\}
 \text{ with }x(N_T(D))\le |D|-1.
\end{gathered}
\end{equation}
Thus admissibility includes failure of every target, in addition to the
anchor constraints.

The cover can be reduced without further search. If $S'\subseteq S$ and
$b'\ge b$, then $x(S)\le b$ implies $x(S')\le b'$ for nonnegative $x$.
Together with coordinatewise inclusion of the root boxes, this identifies
49 cases contained in other cases, leaving 308. Write $u$ for the root upper
bounds. In 29 retained cases, $b+u(S^c)<15$ already contradicts the prescribed
total, so only 279 cases
need further exclusion proofs. The exact map and retained case numbers are
in \path{data/q5-root-reduction.json}; Lean checks all containments and
arithmetic exclusions. The archive retains all 357 original cases.

\begin{lemma}[Hall-budget projection]\label{lem:projection}
Suppose a root has anchor multiplicity $m$, budget $x(S)\le b$, and total
$15$, with the anchor outside $S$. For $0\le \lambda\le b$, put $G=15-b+\lambda$
and suppose $G\ge m$. Every admissible vector $x$ has a vector $y\le x$
with the same anchor multiplicity, total $G$, budget $y(S)\le \lambda$, and a
Hall deficiency for every target. It obeys the original upper bounds,
tightened to $\min(u_i,\lambda)$ on $S$.
\end{lemma}
\begin{proof}
First retain only $\min(x(S),\lambda)$ copies in $S$. This removes at most $b-\lambda$
copies and leaves the anchor untouched, so the remaining total is at least
$G$. Delete further non-anchor copies until the total is $G$; this is possible
because $G\ge m$. Every coordinate decreases, so all upper bounds and Hall
deficiency inequalities are preserved. Nonnegativity and $y(S)\le \lambda$ give
$y_i\le \lambda$ for $i\in S$.
\end{proof}

For our anchored covers, $b=7-m-|K|$, so
\[
 G=m+8+|K|+\lambda.
\]
Thus the reduced collection consists of the $m$ fixed anchor copies and
$8+|K|+\lambda$ other copies; in particular, $G\ge m$ is automatic. Apply the lemma
after selecting and normalising the original root, retaining its bounds.
No new anchor selection is required for the smaller collection.
The resulting 279 projected problems have totals between 9 and 15.
Their parameters and exceptional cases are recorded in
Appendix~\ref{app:verification}.

\section{Interval certificates}\label{sec:certificates}
For a projected problem, admissibility is defined by~\eqref{eq:admissible}
with tightened upper bounds and total $G$. In this section denote its
projected budget by $b$. We certify that no admissible vector exists.

Write a box as $\ell\le x\le u$. For a cut $x(S)\le r$, two necessary
conditions for feasibility in the box are
\begin{equation}\label{eq:possible}
  \ell(S)\le r,\qquad G\le r+u(S^c).
\end{equation}
The first follows from the lower bounds; the second follows from
$G=x(S)+x(S^c)$. When they hold, put
$\sigma=r-\ell(S)$ and $\gamma=r+u(S^c)-G$. Every satisfying vector obeys
\begin{equation}\label{eq:refine}
 x_i\le\min(u_i,\ell_i+\sigma)\quad(i\in S),\qquad
 x_i\ge\max(\ell_i,u_i-\gamma)\quad(i\notin S).
\end{equation}
For example, the upper bound follows by giving all other coordinates of $S$
their lower bounds; the lower bound follows by giving all other coordinates
of $S^c$ their upper bounds.

For a target, retain its cuts satisfying~\eqref{eq:possible}. If none remain,
no vector in the box can fail that target. Otherwise let $I$ and $U$ be the
intersection and union of their neighbourhoods, and let $\sigma_*,\gamma_*$ be the
largest corresponding $\sigma,\gamma$. Then all vectors in the box that fail this
target remain in the refined box
\begin{equation}\label{eq:summary}
 u'_i=\begin{cases}\min(u_i,\ell_i+\sigma_*)&i\in I,\\u_i&i\notin I,\end{cases}
 \qquad
 \ell'_i=\begin{cases}\max(\ell_i,u_i-\gamma_*)&i\notin U,\\\ell_i&i\in U.\end{cases}
\end{equation}
Indeed, a failing vector satisfies one of the retained cuts, and its bounds
in~\eqref{eq:refine} imply~\eqref{eq:summary}. The root budget is treated as
a family containing a single cut. The total $\sum_i x_i=G$ gives
\[
 \ell'_i=\max\!\left(\ell_i,G-\sum_{j\ne i}u_j\right),\qquad
 u'_i=\min\!\left(u_i,G-\sum_{j\ne i}\ell_j\right).
\]
A total-bound leaf closes exactly when $\sum_i\ell_i>G$ or
$\sum_i u_i<G$. A budget leaf closes when
$\ell(S)>b$ or $b+u(S^c)<G$; a target leaf closes when none of its
255 cuts satisfies~\eqref{eq:possible}. The checker thus has 77 closing
reasons: one total test, one budget test and 75 target tests. Refinements
must tighten at least one bound and leave
nonempty coordinate intervals; a split must satisfy $\ell_i\le v<u_i$.

A certificate is a finite tree of refinements and integer splits
$x_i\le v$ or $x_i\ge v+1$. Each leaf is excluded by total bounds or by a
cut family with no surviving cut. Induction on this tree proves that its
root contains no admissible vector. The separate C++ checker reads these
instructions and verifies the deductions; it performs no search and does
not use the search program's matching implementation.

\begin{proposition}[Certificate replay]\label{prop:replay}
The separate checker accepts an exclusion certificate for each of the
279 projected problems.
\end{proposition}
Archive sizes and rejection controls are recorded in
Appendix~\ref{app:verification}.

\begin{proof}[Alternative proof of Theorem~\ref{main-thm:main}]
Lemma~\ref{main-lem:lower} gives $N_5\ge15$. A solution-free collection of fifteen
cubes would satisfy an exported root case by Proposition~\ref{prop:cover},
and would supply a Hall deficiency for every target. Containment and the initial
budget tests leave one of 279 roots. Lemma~\ref{lem:projection} supplies an
admissible vector in its projected problem, which Proposition~\ref{prop:replay}
and the sound interval deductions exclude. Hence $N_5\le15$.
\end{proof}

\section{Verification and reproduction details}\label{app:verification}
\subsection{Finite certificates}
The repeated-variety cases have the following counts,
totalling 244.
\begin{center}
\begin{tabular}{lrrrrrr}
\toprule
Anchor / maximum multiplicity & 2 & 3 & 4 & 5 & 6 & 7\\
\midrule
Adjacent-repeat & 72 & 55 & 31 & 14 & 4 & 1\\
Opposite-repeat & 26 & 21 & 11 & 6 & 2 & 1\\
Total & 98 & 76 & 42 & 20 & 6 & 2\\
\bottomrule
\end{tabular}
\end{center}

Of the 279 remaining roots, 277 admit an exclusion certificate with $\lambda=0$. The other two, original
cases 68 and 69, use $\lambda=1$ and $G=12$. Their $\lambda=0$, $G=11$ problems actually
have solution-free collections, supplied and checked by physical matching in
the projected manifest. The projected totals range from 9 to 15. All 101 roots
whose original budget was zero retain their original certificates.

The projected archive contains 499,829 instructions, 95.69\% fewer than the
original archive. Both archives and the search computations are documented
in the repository.
Every archived header is checked against its exported root case. Eighty
negative controls are rejected, including all 77 possible direct closing
claims against a separate $G=11$ fixture, not a $G=15$ root problem.
This fixture is a known solution-free eleven-cube collection: it satisfies
its total, fixed anchor count, coordinate bounds and zero root budget, and
fails all 75 targets by physical matching. It is \path{data/q5-distinct-known-control.txt}; its regeneration and witness
checks are in \path{scripts/check_five_colour_distinct.py}.

\subsection{Lean implementation}
The development uses Lean~4 and mathlib~\cite{Lean4,mathlib}, both at
\texttt{v4.35.0}. It proves the physical theorems, including the connection
between cube arrangements and the finite certificate problems.
The physical model and palette-independent geometry come from the shared
\texttt{colour-\allowbreak cubes-\allowbreak core} library~\cite{ColourCubes}.

Lean proves completeness and uniqueness of the 75 varieties and checks all
45,000 corner-compatibility bits using the shared corner--rotation theorem.
The lower bound uses 75 deficient corner sets: 72 singletons and three
sets of sizes 5, 6 and 6. Its theorem is
\path{FiveColour.Physical.lower_bound} in \path{FiveColour/LowerBound.lean}. In the formal proof, this replaces the elementary
target classification in Lemma~\ref{main-lem:lower}; catalogue completeness and
rotation transport extend the result to every physical target.

The upper-bound formalisation follows the Hall reduction, symmetry
transport, anchor selection and shared-corner argument. Finite checks
verify capacities, cover representatives and the reduction from 357 roots
to 279. The deletion step is \path{FiveColour.Physical.root_projection}
in \path{FiveColour/RootProjection.lean}. The bridge to the archive, in
\path{FiveColour/ProjectedBridge.lean}, identifies the full admissibility
predicate: coordinate bounds, total, budget and all Hall-cut families.

For efficient replay, Lean proves soundness of a packed implementation of
the interval checker. For multiplicities $0\le x_i\le7$, index the varieties
by $0,\ldots,74$ and put
\[
 P(x)=\sum_{i=0}^{74}x_i B^i,\qquad
 W(S)=\sum_{i\in S}B^{74-i},\qquad B=1024.
\]
Then
\[
 x(S)=\left\lfloor\frac{P(x)W(S)}{B^{74}}\right\rfloor\bmod B.
\]
Every coefficient in the product is at most $75\cdot7<1024$, so no carry
can affect this coefficient extraction. Lean proves the identity, decoding,
coordinate updates and equivalence with the integer inequalities, and
checks all 19,125 neighbourhood weights.

The 279 certificates are split into 4,786 chunks of at most 128
instructions. Composition requires exact boundary-vector agreement and a
proof for every referenced box; the exporter is untrusted. All chunks and
their final assembly have passed kernel checking. The case mapping and
certificate hashes are in \path{data/q5-projected-root-manifest.json}.
The generated module \path{FiveColour/ProjectedChunks/Threshold.lean}
proves the unconditional theorems \path{threshold15} and
\path{sharpThreshold15}. Both lie in namespace
\path{FiveColour.Packed.ProjectedChunks}.

The cuboid proof in Lean follows Lemma~\ref{main-lem:eight-common-missing}.
The cube exports, such as \path{frame20}, specialise the cuboid theorems.
The equivalences in \path{FiveColour/CubeBoxBridge.lean} identify the cube
and equal-side cuboid position models, preserving exposed faces and physical
placements. Cube completion likewise reuses rectangular completion, whose
generic construction is adapted from~\cite{FourToThree}; no theorem of that
paper is used. The auxiliary eleven-cube missing-pair lemma remains in
\path{FiveColour/RingEleven.lean}, with its multiplicity argument proved
without enumerating collections.

The small-support proofs check pair witnesses directly, without reducing
them to five colour orbits, and use deficient corner sets for the two
obstruction constructions. The modules are \path{SmallSupport.lean} and
\path{MixedObstructions.lean} in the \path{FiveColour} directory.

The full cuboid classification is formalised in
\path{FiveColour/BoxThresholds.lean}. In namespace
\path{FiveColour.Physical}, \path{box_frame_threshold_iff} and
\path{box_threshold_iff} give necessary and sufficient pool sizes for
arbitrary side lengths at least two, assuming
\path{FiveColour.Physical.Threshold15}. \path{Solution.lean} applies the
completed corner theorem to discharge this assumption. In namespace
\path{FiveColour.Puzzle}, it exports the unconditional theorems
\path{box_frame_threshold} and \path{complete_box_threshold}, together with
\path{frame20}, \path{frame_threshold_all}, \path{full_cube},
\path{complete_cube_threshold} and \path{support_threshold}.

The rectangular model and completion are in
\path{FiveColour/BoxModel.lean}; compatibility is connected to proper
rotations, and axis assignments are realised by proper rotations.
\path{FiveColour/EightPlusThree.lean} proves the physical prescribed-family
lemma. Its proof uses a failed Hall set to obtain the small cover directly,
as in Lemma~\ref{main-lem:eight-common-missing}; signed projection and small matching
patterns are checked in the kernel. The weighted bounds and upper deductions
are in \path{FiveColour/WideFrames.lean}, \path{FiveColour/RodFrames.lean}
and \path{FiveColour/StripFrames.lean}. The thin-box proofs formalise both
concentration and switching targets, including the equality cases of the
affine geometric bounds. Directional profiles are checked for the two
canonical target varieties; proved colour and rotation transport covers
all targets, including their corner, edge and axis labels. The small
half-mask matching step uses greedy selection after single-mask checks.
\path{FiveColour/BoxLowerBound.lean} proves the
uniform three-variety obstruction for every shape and orientation.
A separate C++ audit reconstructs the geometry and checks these profiles.
No new interval certificates are used. The flow algorithm and its
complexity bound in Section~\ref{main-sec:construction} are mathematical
deductions, not additional Lean claims.

\subsection{Reproduction}
The dependency audits cover 133 declarations listed in \path{Audit.lean},
fourteen extension statements listed in \path{ExtensionsAudit.lean}, and the
final replay theorems. Their dependencies are confined to
\texttt{propext}, \texttt{Classical.choice} and \texttt{Quot.sound}. No \texttt{native\_decide} is used. Reproduce the proof
and audits with:
\begin{verbatim}
make verify-full WORKERS=2
\end{verbatim}
The module guide is \path{FiveColour/README.md}.

A Mathlib-only comparator challenge states the corner problem geometrically,
restating the library's definitions locally. It covers only the corner
problem: the cuboid statements rest on the definitions in
\path{FiveColour/BoxModel.lean}. Its equivalence with the development's
model is proved by
\path{FiveColour.Puzzle.hasSolution_iff} in
\path{FiveColour/ComparatorBridge.lean}. The solution module applies the
completed threshold theorem. A full external comparator pass remains
unconfirmed; it is separate from the completed Lean proof.

\begin{thebibliography}{9}
\raggedright
\bibitem{MacMahon} P.~A.~MacMahon,
\emph{New Mathematical Pastimes}, Cambridge University Press (1921).
\href{https://archive.org/details/newmathematicalp00macmuoft}{Digitised edition}.
\bibitem{BES} E.~Berkove, A.~Eickemeyer and M.~Schulman,
\emph{The 4-Color Cubes Puzzle}, Integers \textbf{18} (2018), A68.
\href{https://doi.org/10.5281/zenodo.10680216}{doi:\nolinkurl{10.5281/zenodo.10680216}}.
\href{https://math.colgate.edu/~integers/s68/s68.pdf}{Journal PDF}.
\bibitem{Hall} P.~Hall, \emph{On Representatives of Subsets},
Journal of the London Mathematical Society \textbf{s1-10}(1) (1935), 26--30.
\href{https://doi.org/10.1112/jlms/s1-10.37.26}{doi:\nolinkurl{10.1112/jlms/s1-10.37.26}}.
\bibitem{EdmondsKarp} J.~Edmonds and R.~M.~Karp,
\emph{Theoretical Improvements in Algorithmic Efficiency for Network Flow Problems},
Journal of the ACM \textbf{19}(2) (1972), 248--264.
\href{https://doi.org/10.1145/321694.321699}{doi:\nolinkurl{10.1145/321694.321699}}.
\bibitem{FourToThree} J.~Davies,
\emph{Three-colour solutions from four-colour cubes}, manuscript (2026).
\href{https://github.com/jasondavies/four-to-three-colour-cubes}{Accompanying repository}.
\bibitem{ColourCubes} J.~Davies,
\texttt{colour-\allowbreak cubes-\allowbreak core}: palette-independent Lean definitions and proofs
for face-coloured cubes, Lean library (2026).
Release \href{https://github.com/jasondavies/colour-cubes-core/tree/v1.0.0}{\texttt{v1.0.0}},
\url{https://github.com/jasondavies/colour-cubes-core}.
\bibitem{Lean4} L.~de Moura and S.~Ullrich,
\emph{The Lean 4 Theorem Prover and Programming Language},
in Automated Deduction---CADE~28, LNCS \textbf{12699}, Springer (2021), 625--635.
\href{https://doi.org/10.1007/978-3-030-79876-5_37}{doi:\nolinkurl{10.1007/978-3-030-79876-5_37}}.
\bibitem{mathlib} The mathlib Community, \emph{The Lean Mathematical Library},
in Proceedings of CPP~2020, ACM (2020), 367--381.
\href{https://doi.org/10.1145/3372885.3373824}{doi:\nolinkurl{10.1145/3372885.3373824}}.
\href{https://arxiv.org/abs/1910.09336}{arXiv:1910.09336}.
\end{thebibliography}
\end{document}
